\chapter{Logic}
\label{ch:logic}

Logic is the starting point, and it can't be constructed from anything earlier.
It is the set of rules we agree to reason by. Everything after this chapter is
built using these rules, and only these rules.

We use \emph{classical logic}: every statement is either true or false, never
both and never neither.

\section{Propositions}

\begin{definition}[Proposition]
\label{def:proposition}
A \emph{proposition} is a statement that is either true or false. Its \emph{truth
value} is $\T$ (true) or $\F$ (false).
\register{logic}{$\T$, $\F$}{the truth values true and false}
\end{definition}

For example, \enquote{every set is a subset of itself} is a proposition.
\enquote{$x$ is a subset of $y$} is not, since its truth depends on $x$ and $y$.
That kind of statement is handled in \cref{sec:quantifiers}. We write propositions
as $P, Q, R, \dots$

\section{Connectives}

New propositions are formed from old ones with \emph{connectives}. A connective is
defined by its truth table: the truth value of the new proposition for every
possible truth value of its parts.

\begin{definition}[Connectives]
\label{def:connectives}
\[
\begin{array}{cc|ccccc}
  P & Q & \lnot P & P \land Q & P \lor Q & P \Rightarrow Q & P \Leftrightarrow Q \\
  \hline
  \T & \T & \F & \T & \T & \T & \T \\
  \T & \F & \F & \F & \T & \F & \F \\
  \F & \T & \T & \F & \T & \T & \F \\
  \F & \F & \T & \F & \F & \T & \T
\end{array}
\]
\begin{itemize}
  \item $\lnot P$ (\emph{not} $P$, negation);
  \item $P \land Q$ ($P$ \emph{and} $Q$, conjunction);
  \item $P \lor Q$ ($P$ \emph{or} $Q$, disjunction). This is the inclusive
    \enquote{or}: it is true when both are true;
  \item $P \Rightarrow Q$ (\emph{if} $P$ \emph{then} $Q$, implication). $P$ is the
    \emph{hypothesis} and $Q$ the \emph{conclusion};
  \item $P \Leftrightarrow Q$ ($P$ \emph{if and only if} $Q$, equivalence).
\end{itemize}
\register{logic}{$\lnot P$}{not $P$}
\register{logic}{$P \land Q$}{$P$ and $Q$}
\register{logic}{$P \lor Q$}{$P$ or $Q$ (inclusive)}
\register{logic}{$P \Rightarrow Q$}{if $P$ then $Q$}
\register{logic}{$P \Leftrightarrow Q$}{$P$ if and only if $Q$}
\end{definition}

\begin{remark}[Vacuous truth]
\label{rem:vacuous}
$P \Rightarrow Q$ only claims: \emph{whenever $P$ is true, $Q$ is true as well}.
It makes no claim when $P$ is false, so it counts as true then (rows 3 and 4).
This is called \emph{vacuous truth}: \enquote{every element of the empty
collection is red} is true, because no element can contradict it.
\end{remark}

\begin{definition}[Tautology, contradiction, equivalence]
\label{def:tautology}
A compound proposition is a \emph{tautology} if it is true for every assignment of
truth values to its parts, and a \emph{contradiction} if it is false for every
assignment. Two compound propositions $A$ and $B$ are \emph{logically
equivalent}, written $A \equiv B$, if $A \Leftrightarrow B$ is a tautology, so
they have the same truth table.
\register{logic}{$A \equiv B$}{$A$ and $B$ are logically equivalent}
\end{definition}

\begin{theorem}[Basic laws]
\label{thm:logic-laws}
For all propositions $P, Q, R$:
\begin{center}
\begin{tabular}{@{}ll@{}}
  \toprule
  Law & Statement \\
  \midrule
  Double negation & $\lnot\lnot P \equiv P$ \\
  Excluded middle & $P \lor \lnot P$ is a tautology \\
  Non-contradiction & $P \land \lnot P$ is a contradiction \\
  Commutativity & $P \land Q \equiv Q \land P$, \quad $P \lor Q \equiv Q \lor P$ \\
  Associativity & $(P \land Q) \land R \equiv P \land (Q \land R)$, \quad same for $\lor$ \\
  Distributivity & $P \land (Q \lor R) \equiv (P \land Q) \lor (P \land R)$ \\
    & $P \lor (Q \land R) \equiv (P \lor Q) \land (P \lor R)$ \\
  De Morgan & $\lnot(P \land Q) \equiv \lnot P \lor \lnot Q$, \quad
    $\lnot(P \lor Q) \equiv \lnot P \land \lnot Q$ \\
  Implication as \enquote{or} & $P \Rightarrow Q \equiv \lnot P \lor Q$ \\
  Contrapositive & $P \Rightarrow Q \equiv \lnot Q \Rightarrow \lnot P$ \\
  Negated implication & $\lnot(P \Rightarrow Q) \equiv P \land \lnot Q$ \\
  Equivalence & $P \Leftrightarrow Q \equiv (P \Rightarrow Q) \land (Q \Rightarrow P)$ \\
  \bottomrule
\end{tabular}
\end{center}
\end{theorem}

\begin{proof}
Each law is checked by writing out the truth table of both sides for all $2$, $4$
or $8$ assignments and comparing them row by row. We show one case, De Morgan's
first law:
\[
\begin{array}{cc|ccc|ccc}
  P & Q & P \land Q & \lnot(P \land Q) & & \lnot P & \lnot Q & \lnot P \lor \lnot Q \\
  \hline
  \T & \T & \T & \mathbf{F} & & \F & \F & \mathbf{F} \\
  \T & \F & \F & \mathbf{T} & & \F & \T & \mathbf{T} \\
  \F & \T & \F & \mathbf{T} & & \T & \F & \mathbf{T} \\
  \F & \F & \F & \mathbf{T} & & \T & \T & \mathbf{T}
\end{array}
\]
The bold columns agree. The other laws are checked the same way.
\end{proof}

\begin{remark}
The \emph{converse} $Q \Rightarrow P$ is \emph{not} equivalent to $P \Rightarrow
Q$. Take $P$ false and $Q$ true: then $P \Rightarrow Q$ is true and $Q \Rightarrow
P$ is false.
\end{remark}

\section{Valid inference}

\begin{definition}[Valid inference]
\label{def:valid}
An inference from premises $A_1, \dots, A_n$ to a conclusion $B$ is \emph{valid}
if $(A_1 \land \dots \land A_n) \Rightarrow B$ is a tautology. In that case,
whenever all premises are true, the conclusion is true too.
\end{definition}

\begin{theorem}[Basic inference rules]
\label{thm:inference}
The following inferences are valid:
\begin{center}
\begin{tabular}{@{}lll@{}}
  \toprule
  Rule & Premises & Conclusion \\
  \midrule
  Modus ponens & $P$, \ $P \Rightarrow Q$ & $Q$ \\
  Modus tollens & $\lnot Q$, \ $P \Rightarrow Q$ & $\lnot P$ \\
  Chaining & $P \Rightarrow Q$, \ $Q \Rightarrow R$ & $P \Rightarrow R$ \\
  Cases & $P \lor Q$, \ $P \Rightarrow R$, \ $Q \Rightarrow R$ & $R$ \\
  Contradiction & $\lnot P \Rightarrow (Q \land \lnot Q)$ & $P$ \\
  \bottomrule
\end{tabular}
\end{center}
\end{theorem}

\begin{proof}
Again by truth tables. For modus ponens: $(P \land (P \Rightarrow Q)) \Rightarrow
Q$ could only be false if $Q$ is false and $P \land (P \Rightarrow Q)$ is true.
Then $P$ is true, and $P \Rightarrow Q$ with $P$ true and $Q$ false is false. That
is a contradiction, so the implication is never false. The others are similar.
\end{proof}

\section{Predicates and quantifiers}
\label{sec:quantifiers}

\begin{definition}[Predicate]
\label{def:predicate}
A \emph{predicate} is a statement containing a \emph{variable} $x$, which becomes
a proposition once a specific object is put in for $x$. We write it $P(x)$, and
$P(a)$ for the proposition obtained by putting in $a$. Predicates can contain
several variables: $P(x, y)$.
\register{logic}{$P(x)$}{a statement in which the variable $x$ occurs; not a map
value (\cref{rule:predicates})}
\end{definition}

\begin{definition}[Quantifiers]
\label{def:quantifiers}
\leavevmode
\begin{itemize}
  \item $\forall x\; P(x)$ (\enquote{for all $x$, $P(x)$}) is true if and only if
    $P(a)$ is true for every object $a$.
  \item $\exists x\; P(x)$ (\enquote{there exists $x$ with $P(x)$}) is true if and
    only if $P(a)$ is true for at least one object $a$.
  \item $\exists! x\; P(x)$ (\enquote{there is exactly one $x$ with $P(x)$})
    abbreviates
    \[
      \exists x\; \bigl( P(x) \land \forall y\; (P(y) \Rightarrow y = x) \bigr).
    \]
\end{itemize}
The variable after a quantifier is \emph{bound}: $\forall x\, P(x)$ and $\forall
y\, P(y)$ say the same thing. A variable that isn't bound is \emph{free}.
\register{logic}{$\forall x\; P(x)$}{for all $x$, $P(x)$}
\register{logic}{$\exists x\; P(x)$}{there is an $x$ with $P(x)$}
\register{logic}{$\exists! x\; P(x)$}{there is exactly one $x$ with $P(x)$}
\end{definition}

\begin{definition}[Equality]
\label{def:equality}
Our logic includes the symbol $=$. The statement $a = b$ means $a$ and $b$ are the
same object. It satisfies $a = a$, and if $a = b$ then anything true of $a$ is
true of $b$ (\emph{substitution}). We write $a \neq b$ for $\lnot (a = b)$.
\register{logic}{$a = b$}{$a$ and $b$ are the same object}
\register{logic}{$a \neq b$}{$\lnot (a = b)$}
\end{definition}

From the two rules of equality: $a = b \Rightarrow b = a$, and $(a = b \land b = c)
\Rightarrow a = c$.

\begin{definition}[Defining new notation]
\label{def:defining}
$A \defiff B$ means that the new statement $A$ is defined to mean the statement
$B$. $a \coloneqq t$ means that the new name $a$ is defined to be the object $t$.
\register{logic}{$A \defiff B$}{the statement $A$ is defined to mean $B$}
\register{logic}{$a \coloneqq t$}{the name $a$ is defined to be $t$}
\end{definition}

\begin{theorem}[Negating quantifiers]
\label{thm:negate-quantifiers}
\[
  \lnot \forall x\; P(x) \;\equiv\; \exists x\; \lnot P(x), \qquad
  \lnot \exists x\; P(x) \;\equiv\; \forall x\; \lnot P(x).
\]
\end{theorem}

\begin{proof}
$\forall x\, P(x)$ is false exactly when some object $a$ makes $P(a)$ false, which
means $\lnot P(a)$ is true for some $a$. That is the first equivalence. For the
second, apply the first to the predicate $\lnot P$ and use double negation:
$\lnot \forall x\, \lnot P(x) \equiv \exists x\, \lnot\lnot P(x) \equiv \exists
x\, P(x)$. Negating both sides gives the claim.
\end{proof}

For example, the negation of \enquote{every set has an element} is \enquote{some
set has no element}.

\begin{remark}[Order of quantifiers]
\label{rem:quantifier-order}
Quantifiers of the same kind can be swapped,
\[
  \forall x\, \forall y\; P(x,y) \;\equiv\; \forall y\, \forall x\; P(x,y),
\]
and the same holds for $\exists$. Mixed quantifiers cannot. In general
\[
  \exists y\, \forall x\; P(x,y) \;\Rightarrow\; \forall x\, \exists y\; P(x,y)
\]
holds, but the converse doesn't. \enquote{For every number there is a bigger one}
is true for the natural numbers. \enquote{There is a number bigger than every
number} is false.
\end{remark}

\begin{definition}[Bounded quantifiers]
\label{def:bounded}
\[
  \forall x \in A\; P(x) \defiff \forall x\, (x \in A \Rightarrow P(x)), \qquad
  \exists x \in A\; P(x) \defiff \exists x\, (x \in A \land P(x)),
\]
and $\exists! x \in A\; P(x)$ means there is exactly one $x$ with $x \in A$ and
$P(x)$. Several variables are bound at once by $\forall x, y \in A$.
\register{logic}{$\forall x \in A\; P(x)$}{for all $x$ in $A$, $P(x)$}
\register{logic}{$\exists x \in A\; P(x)$}{there is an $x$ in $A$ with $P(x)$}
\end{definition}

By \cref{thm:negate-quantifiers}, $\lnot \forall x \in A\, P(x) \equiv \exists x
\in A\, \lnot P(x)$, and similarly for $\exists$. (The relation $\in$ is introduced
in \cref{ch:sets}.)

\section{Proof methods}
\label{sec:proof-methods}

A \emph{proof} is a finite chain of statements. Each one is an axiom, a
definition, a result already proved, an assumption made for the proof, or follows
from earlier statements by a valid inference. These are the standard patterns:

\begin{description}[style=nextline]
  \item[Direct proof] of $P \Rightarrow Q$: assume $P$, derive $Q$.
  \item[Contrapositive] of $P \Rightarrow Q$: assume $\lnot Q$, derive $\lnot P$.
    Valid by \cref{thm:logic-laws}.
  \item[Contradiction] for $P$: assume $\lnot P$, derive some statement together
    with its negation. Valid by \cref{thm:inference}.
  \item[Cases:] if $P_1 \lor \dots \lor P_n$ holds and $R$ follows from each
    $P_i$, then $R$ holds.
  \item[Equivalence] $P \Leftrightarrow Q$: prove $P \Rightarrow Q$ and $Q
    \Rightarrow P$ separately.
  \item[Universal statement] $\forall x\, P(x)$: let $x$ be \emph{arbitrary},
    assuming nothing about it, and prove $P(x)$.
  \item[Existence] $\exists x\, P(x)$: name a specific object $a$ and prove
    $P(a)$. Alternatively, prove it by contradiction.
  \item[Uniqueness] in $\exists! x\, P(x)$: prove existence, then assume $P(x)$ and
    $P(y)$ and derive $x = y$.
\end{description}

\emph{Induction} is added in \cref{ch:naturals}. There it is a theorem, proved
from the axioms, not a new rule of logic.

\begin{exercises}
  \item Show by truth table that $(P \Rightarrow Q) \land (P \Rightarrow R) \equiv
    P \Rightarrow (Q \land R)$.
  \item Is $(P \Rightarrow Q) \Rightarrow R$ equivalent to $P \Rightarrow (Q
    \Rightarrow R)$? Prove it or give a counterexample.
  \item Write the negation of $\forall \varepsilon\, \exists \delta\, \forall x\;
    (P(x, \delta) \Rightarrow Q(x, \varepsilon))$ without a $\lnot$ in front of a
    quantifier. This shape comes back in calculus.
\end{exercises}
